\begin{abstract} 
Low birth weight (LBW) remains a critical public-health indicator, linked strongly with higher neonatal mortality, developmental delays, and lifelong chronic diseases. Using the 2021 U.S. Natality dataset (> 3 million births), this thesis develops a Bayesian, tree-based, nonparametric framework that models the full birth-weight distribution and quantifies LBW risk. 

The raw dataset is condensed into 672 mutually exclusive classes defined by seven maternal-infant predictors and 10 and 11 birth-weight categories. The former is comprised of 10\% LBW quantile categories while the latter adds an additional aggregated normal birth-weight (NBW) category. These different models are grown to contrast and investigate how variable selection differs based on the restriction of the dataset. These models are Classification and Regression Trees (CART) using the marginal Dirichlet-Multinomial likelihood as the splitting criterion. This criterion is equipped to handle sparse observations, with the Dirichlet hyperparameters informed by previous quantiles from the 2023 dataset to avoid "double dipping."

To ensure stable growth of the trees, a two-tier parametric bootstrap resampling technique is employed, yielding a 10,000 tree ensemble. This ensemble provides highly stable prediction estimates. Maternal race, smoking status, and marital status consistently drive the initial LBW risk stratification, identifying black, smoking, unmarried mothers among the highest-risk subgroups. When the analysis is restricted to LBW births only, infant gender and maternal age supersede smoking and marital status as key discriminators, revealing finer biologically-driven risk factors. Stable and informative mean ensemble estimates are obtained with narrow 95\% percentile intervals.

% The full and LBW-only models are grown to contrast and investigate how variable selection is altered based on the restriction the dataset. The models are  Classification and Regression Trees (CART) using the marginal Dirichlet-Multinomial likelihood as the splitting criterion. This criterion is equipped to handle sparse observations, with the Dirichlet hyperparameters informed by previous quantiles from the 2020 dataset to avoid "double dipping."

% Employing a two-tier parametric bootstrap resampling technique, a 10,000 tree ensemble is grown yielding highly stable prediction estimates. Maternal race, smoking status, and marital status consistently drive the initial LBW risk stratification, identifying black, smoking, unmarried mothers among the highest-risk subgroups. When the analysis is restricted to LBW births only, infant gender and maternal age supersede smoking and marital status as key discriminators, revealing finer biological gradients of risk. Stable and informative mean ensemble estimates are obtained with narrow 95\% percentile intervals.

The resulting modeling framework combines the interpretability of decision trees with a custom quasi-Bayesian splitting criterion, yielding delivering actionable, clinically relevant insights for targeting maternal-health interventions among the most vulnerable subpopulations.
\end{abstract}